\section{Related Work}
\label{sec:related-work}

\subsection{LLM-Based Query Expansion}

LLM-based query expansion supplements short queries with document-like
language that may be absent from the original formulation. HyDE encodes a
generated hypothetical document for dense retrieval, whereas Query2doc appends
a pseudo-document to the query for sparse and dense retrieval
\citep{gao-etal-2023-precise,wang-etal-2023-query2doc}. Later methods generate
different forms of evidence or sample multiple references to cover alternative
expressions of the same information need
\citep{jagerman-etal-2023-query-expansion,zhang-etal-2024-exploring-best}.
Although generated text can supply useful lexical and semantic cues, its effect
varies across queries, datasets, and retrievers, and irrelevant content may
reduce retrieval effectiveness \citep{weller-etal-2024-generative}.

\subsection{Evaluating Expansion under Hybrid Fusion}

Query-expansion methods commonly evaluate retrieval using a fixed top-$L$
prefix from each ranking
\citep{zhang-etal-2024-exploring-best,liu-zhang-2025-exp4fuse}. In hybrid
retrieval, $L$ also determines which cross-channel contributions enter fusion,
making the measured expansion effect conditional on the selected cutoff.
Following EAHR, we evaluate complete-list retrieval effectiveness separately from
channel access depth \citep{zhang2026exact}, and ask whether query expansion can
improve the former while reducing the latter.

\subsection{Channel-Aware Integration}

Generated content can enter retrieval through query text, dense
representations, term weights, or rank fusion. Query2doc modifies the query
text, MuGI pools representations from multiple references, Word2Passage
reweights generated terms, and Exp4Fuse combines original and expanded
rankings
\citep{wang-etal-2023-query2doc,zhang-etal-2024-exploring-best,choi-etal-2025-word2passage,liu-zhang-2025-exp4fuse}.
DVCQR further constructs separate lexical and semantic views for sparse and
dense retrieval \citep{li-etal-2026-dvcqr}. These methods differ in where
expansion is integrated; DESA instead focuses on how the same generated passages should
be constrained differently across the two channels.

Recent work also combines heterogeneous retrievers per query: MoR weights
sparse and dense retrievers, while QuDAR adaptively fuses four rankings formed
by original and expanded queries under both retriever types
\citep{kalra-etal-2025-mor,kim-etal-2026-qudar}. DESA instead modifies the dense
and sparse rankings before fusion: generated passages expand dense semantic
coverage while sparse retrieval remains anchored to the original query's
lexical support.
